% =====================================================================
%  INTERNSHIPS & INDUSTRIAL TRAINING — CV section
%  Drop-in replacement for the existing "Internship" block.
%
%  NOTE ON THE HEADING: an industrial attachment is not an internship.
%  Filing it under "Internship" invites a reviewer to read it as paid
%  work you performed. "Internships & Industrial Training" covers both
%  honestly and costs you nothing.
%
%  NOTE ON VERBS: this was a supervised plant study, not a work
%  placement. Every verb below is one you can defend in an interview -
%  studied, traced, analysed, examined, documented, reviewed. Do not
%  upgrade them to "designed", "operated" or "implemented".
%
%  Pick 3-4 APSCL bullets for a 2-page CV; use all six only if the
%  target lab is power systems and the space exists.
% =====================================================================

\section{Internships \& Industrial Training}

\textbf{Industrial Attachment --- Ashuganj Power Station Company Limited (APSCL)} \hfill Nov 2025\\
Brahmanbaria, Bangladesh \quad|\quad Electrical Energy and Power System
\begin{itemize}
  \item Studied a \textbf{1,647~MW} combined-cycle gas generation complex end to end under APSCL executive engineers, tracing the full conversion chain from gas intake through compressor, combustion, gas turbine, HRSG, steam turbine and generator to the 400~kV bus
  \item Analysed how heat recovery steam generation lifts plant efficiency from under 35\% to \textbf{$\approx$58\%}, mapping the three-pressure HRSG train (120/31/3~bar) across its economiser, evaporator, superheater and reheat stages
  \item Examined parallel-generator control in operation --- droop characteristics, \textbf{FGMO} isochronous load sharing with master/slave governors and a load sharing controller, and governor-induced load hunting
  \item Compared 400~kV \textbf{AIS and GIS} switchgear and \textbf{Double Bus Double Breaker} topology, weighing $n{-}1$ redundancy against capital cost, footprint and SF\textsubscript{6} leak management
  \item Documented the \textbf{black-start} sequence from scheduled shutdown: SFC-driven barring at 120~rpm, generator handover at 60--63\% of rated speed, and staged diverter-damper opening to avoid HRSG thermal shock
  \item Reviewed \textbf{30+} generator, transformer and busbar protection schemes alongside DCS monitoring, fail-proof setpoints, majority-rule transmitter voting and cascaded valve actuation
\end{itemize}

\vspace{5pt}

\textbf{VLSI Engineering Intern, Siliconova Ltd.} \hfill Nov 2025\\
<<City>>, Bangladesh
\begin{itemize}
  \item <<Verb>> <<what you personally did on the RTL or verification side>>, <<result>>
  \item <<Verb>> <<what you produced in DFT, physical design or layout>>, <<result>>
\end{itemize}
\href{https://drive.google.com/file/d/16FEPvnWIKaB9XMFpjI0BC0aNkjKeaeZY/view}{\faCertificate\ Certificate}

% -----------------------------------------------------------------
%  The Siliconova entry currently reads:
%
%    "Gained hands-on exposure to end-to-end VLSI design flow covering
%     RTL, verification, DFT, physical design, analog design, and
%     layout techniques."
%
%  Two problems. "Gained hands-on exposure" is passive and tells a
%  reader nothing about what you did. And six sub-fields in one line
%  reads as a syllabus, not an experience - no one does RTL, DFT,
%  physical design AND analog layout meaningfully in one month.
%
%  Pick the one or two areas you actually touched and say what came
%  out of it. If the honest answer is that it was a taught programme
%  rather than project work, say so plainly and move the entry under
%  Industrial Training with the same verbs used above. A short honest
%  entry survives an interview; a broad vague one does not.
% -----------------------------------------------------------------
